\subsection{Connected-region classification}\label{app:connectedproof}
We retain the loopless multigraph conventions of
Appendix~\ref{app:notation}; reciprocal arcs are distinct undirected
edges when defining atoms and connectivity.

\begin{lemma}[localization of a separating bond]\label{lem:bondlocal}
If $G[U]$ and $G[V\setminus U]$ are connected, all edges of
$\delta(U)$ lie in one cycle atom (or its singleton bridge).
If $s\in U$ and $t\notin U$, that atom lies on $\mathcal P(s,t)$.
\end{lemma}
\begin{pf}
For two distinct crossing edges, connect their endpoints by a simple
path inside $U$ and by another inside its complement. The two paths
have disjoint vertex sets; together with the crossing edges they form
a simple cycle, including the possible two-edge cycle. Hence all
crossing edges share a circuit class. A singleton crossing edge is a
bridge. A simple $s$--$t$ path crosses this boundary, so it meets its
atom; every atom met by such a path lies on $\mathcal P(s,t)$.
\end{pf}

\begin{lemma}[connected prefixes in an atom]\label{lem:prefixatom}
Let $A$ be a non-bridge atom and let $a\ne b$ be two of its vertices.
Its vertices admit an ordering $v_1=a,\ldots,v_k=b$ such that every
proper nonempty prefix and its suffix induce connected subgraphs.
The assertion includes two-vertex parallel-edge atoms.
\end{lemma}
\begin{pf}
For two vertices the assertion is immediate. Otherwise add an $a$--$b$
edge temporarily if needed. The connected-prefix property is the classical
$st$-numbering property \citep{even1976st}; we give a construction here.
A biconnected graph with this edge admits
an open-ear construction starting from a cycle containing that edge.
For completeness, this construction follows directly from vertex
connectivity. If vertices remain outside the current subgraph, any
component of those vertices has at least two distinct attachment
vertices in the current subgraph: a unique attachment would be an
articulation of the original graph. A path through the component
between two attachments supplies the next ear. Repeat until all
vertices have entered. Edges left over are irrelevant to the ordering.
The initial cycle exists because the added edge is not a bridge.

Order the other arc of the initial cycle from $a$ to $b$. For each
subsequent ear, let $u$ precede $v$ among its two attachment vertices;
insert its internal vertices, in their order from $u$ to $v$,
immediately after $u$. Every new internal vertex has an earlier and a
later neighbor. Existing vertices keep these neighbors. The first and
last vertices remain $a,b$. Every vertex except $a$ therefore has an
earlier neighbor and every vertex except $b$ a later neighbor. Repeated
steps along these neighbors connect every prefix to $a$ and every
suffix to $b$. Removing the temporary edge does not affect this
argument: it is never internal to a proper prefix or suffix.
\end{pf}

\begin{pf}[Theorem~\ref{thm:connectedclassification}]
For sufficiency, Lemma~\ref{lem:bondlocal} places every crossing edge
in one traversed atom. Its weight is $\alpha$, so
$-\mathbf1_U^\top BDx=-\alpha\mathbf1_U^\top Bx
=-\alpha\mathbf1_U^\top(e_t-e_s)=\alpha$.
No other edge enters this boundary sum.

For necessity, fix a traversed atom $A$. Delete its edges but keep its
vertices. The remaining graph has one connected component $T_v$ for
each $v\in V(A)$, containing that vertex and no other vertex of $A$.
Indeed an outside path between two vertices of $A$ would join their
blocks in a cycle, contrary to the block-cut tree. Connectivity of $G$
shows these components partition $V$. There are distinct vertices
$a,b$ with $s\in T_a$ and $t\in T_b$; these are the entry and exit of
the OD block path in $A$. For a bridge, take the cut consisting of
its $s$-side component. Since its unit conservative crossing response
is fixed, the assumed account immediately gives that bridge's weight
$\alpha$.

For a non-bridge atom, choose the ordering of
Lemma~\ref{lem:prefixatom}. Set
$U_i=\bigcup_{\ell\le i}T_{v_\ell}$ for $1\le i<k$.
The bags and the prefix/suffix connectivity show
$U_i\in\mathcal U^{\mathrm{conn}}_{st}$. Its crossing edges are exactly
those of the corresponding prefix in $A$.
Now choose any signed local response $y$ with
$B_Ay=e_b-e_a$, where $B_A$ is the incidence matrix on $A$.
Extend it by a signed unit path from $s$ to $a$ inside $T_a$ and from
$b$ to $t$ inside $T_b$, with zero on other edges. The result satisfies
$Bx=\mu$. Thus the assumption gives
\[
-\sum_{\ell\le i}(B_AD_Ay)_{v_\ell}=\alpha
\quad(1\le i<k).
\]
The first prefix gives $(B_AD_Ay)_a=-\alpha$, differences of consecutive
prefixes give zero at every interior vertex, and zero total incidence
sum gives $(B_AD_Ay)_b=\alpha$. Consequently
$B_AD_Ay=\alpha(e_b-e_a)$ for every local conservative $y$.
The single-OD classification (Theorem~\ref{thm:t3cpmain}), applied to the
one-atom graph $A$, forces every edge in $A$ to have weight $\alpha$.
This holds for every traversed atom. The sufficiency argument already
shows that no off-path weight is constrained.
\end{pf}

\begin{pf}[Corollary~\ref{cor:connectedmulti}]
Apply Theorem~\ref{thm:connectedclassification} at the common $w$ to
each query. Each queried path requires one weight on all its edges;
intersecting paths propagate equality transitively. Conversely these
transitive equalities satisfy every single-query condition. Atoms not
on any queried path enter none of these equalities.
\end{pf}

\subsection{Finite-change supervision and joint-path implication}\label{app:pairproof}
\begin{pf}[Proposition~\ref{prop:pairboundary}]
Write $r\in\{L,S\}$ and $c^r(q)=g(f^r(q))$ componentwise. Along a
single-coordinate segment $q(u)=q^-+u(q^+-q^-)$, the chain rule gives
\[
\frac{d}{du}c^r(q(u))=(q_j^+-q_j^-)
D^r(q(u))J_j^r(q(u)).
\]
Integration gives the attribution column exactly as
$\widetilde C^r_{:,j}=c^r(q^+)-c^r(q^-)=z_j^r$.
Corollary~\ref{cor:boundaryerror} applied to their difference proves
the stated maximum-error identity. No reference input derivative is
needed to obtain this label.

For the local-to-joint implication set
$A(q)=B[c^L(q)-c^S(q)]$. Equality of every separating boundary account
of a sufficiently small single-coordinate change implies
$\mathcal K_j(c^L(q+\epsilon e_j)-c^L(q)
-c^S(q+\epsilon e_j)+c^S(q))=0$.
The seminorm's kernel is $\ker B$: its expression forces all internal
entries of $Bh$ to vanish and its two endpoint entries to agree; their
zero total then forces both to vanish. Hence
$A(q+\epsilon e_j)=A(q)$. Differentiability and division by $\epsilon$
yield
\[
\partial_j A(q)=B[D^L(q)J_j^L(q)-D^S(q)J_j^S(q)]=0.
\]
On any valid joint attribution path, integration of this identity gives
zero divergence of the corresponding attribution-column error, and
therefore matching accounts on every separating boundary. This needs
the local matching premise throughout the domain containing the path;
a finite training sample supplies no such premise.
\end{pf}
